We have investigated what remains to be designed once a variational quantum learning model has been constrained to respect a data symmetry. Equivariance is a powerful inductive bias, but it is not a complete ansatz prescription. It tells the model which inputs to treat as equivalent and how to share parameters across symmetry-related parts of the data, but it leaves open how trainable capacity should be assigned to the structures of the task. Our results show that this allocation is a central part of quantum model selection.

Tic-Tac-Toe exposes this point because its symmetries and decisive winning triples are known exactly. Full $D_4$ is the natural symmetry of the board, and enforcing it improves the standard edge ansatz over an unconstrained circuit. At the same time, the subgroup sweep shows that equivariance is not all-or-nothing: in our setting, $C_4$ sharing recovers almost the same generalization behavior as full $D_4$. Thus, the strength of symmetry enforcement can itself be treated as an architectural hyperparameter. Partial symmetry can preserve most of the useful bias while avoiding some of the rigidity introduced by the full group. The choice also affects trainability: symmetry restrictions can speed up training and mitigate barren plateaus~\cite{sauvage2024building,cerezo2023simulability}, yet strongly constrained circuits may approach classically simulable regimes~\cite{cerezo2023simulability}.

The larger effect comes from the choice of interactions. The edge ansatz respects neighboring board fields, but the label is generated by complete rows, columns, and diagonals. Adding orbit-shared $ZZZ$ and $\mathrm{CCR}_Z$ gates on these winning triples keeps the model equivariant while placing trainable interactions directly on the motifs that define the task. The random-sharing, ablation, and parameter-count controls show that the improvement is not explained by fewer parameters or by adding gates indiscriminately. It comes from aligning the equivariant circuit with the geometry of the target function.

Our contribution is therefore a refinement of the symmetry program for variational quantum learning. A practical recipe is to represent the data symmetry, choose the subgroup to enforce, and place trainable interactions on label-relevant motifs. Tic-Tac-Toe is small and exactly enumerable, but precisely for this reason it exposes this principle without confounding factors. Future work should study softer versions of this recipe, such as regularized deviations or horizontal-gate constructions that interpolate between strict equivariance and fully free circuits~\cite{wiersema2025horizontal}, compare winning-line gates with parameter-matched gates on non-winning triples to isolate motif alignment, and apply motif-aligned equivariant ans\"atze to larger tasks where the relevant structures must be inferred from domain knowledge.
